\documentclass{article}
\usepackage[utf8]{inputenc}
\usepackage[T1]{fontenc}
\usepackage[a4paper,total={6in,8in}]{geometry}
\usepackage{amsmath,amssymb,amsfonts}
\usepackage{graphicx}
\usepackage{hyperref}
\usepackage{enumitem}
\usepackage{listings}
\usepackage{xcolor}
\DeclareMathOperator{\Var}{Var}
\DeclareMathOperator{\Cov}{Cov}
\providecommand{\tightlist}{\setlength{\itemsep}{0pt}\setlength{\parskip}{0pt}}
\lstset{basicstyle=\ttfamily\scriptsize,breaklines=true,columns=fullflexible,
  keepspaces=true,showstringspaces=false,frame=single}
\hypersetup{colorlinks=true,urlcolor=blue,linkcolor=blue}
\setlength{\emergencystretch}{2em}

\title{Math 126 Problem 6: Set Size Estimate Accuracy}
\author{} % Put your name here.
\date{Spring 2027}
\begin{document}
\maketitle

\hypertarget{submission-instructions}{%
\subsection*{Submission Instructions}\label{submission-instructions}}

Please upload your work to Gradescope by \textbf{11:59pm on Monday, February 15, 2027 (Pacific time).}

\begin{itemize}
\item
  You are encouraged to discuss ideas and work with your classmates. However, you \emph{must acknowledge} your collaborators at the top of each solution on which you collaborated with others and you \emph{must write} your solutions independently.
\item
  Your solutions to theory questions must be written legibly, or typeset in LaTeX or markdown.
\item
  I recommend that you write your solutions in a Python notebook using Google Colab.
\item
  You should submit your solutions as a \emph{single PDF} via the assignment on Gradescope.
\end{itemize}

\textbf{Grading:} The point of the problem set is for \emph{you} to learn. To this end, I hope to disincentivize the use of LLMs by \emph{not} grading your work for correctness. Instead, you will grade your own work by comparing it to my solutions. This self-grade is due the Friday \emph{after} the problem set is due, also on Gradescope.

\noindent Original notebook: \href{https://colab.research.google.com/drive/1oTX14kBIP2e7m95LHlsqStddT8FTAzKY?usp=sharing}{Google Colab}.

% Source notebook cell 2
\hypertarget{set-size-estimate-accuracy}{%
\subsection*{Set Size Estimate Accuracy}\label{set-size-estimate-accuracy}}

Recall the set size estimation problem, where we draw \(m\) uniform samples from a set of size \(n\). Let \(D\) be the number of duplicates in these samples that we observed. On the notes, we claimed that:

If we make \(m \geq c \frac{\sqrt{n}}{\epsilon}\) samples for a particular constant \(c\), then the estimate \(\widehat{n} = \frac{m(m-1)}{2D}\) satisfies \((1-\epsilon) n \leq \widehat{n} \leq (1+\epsilon) n\) with probability \(9/10\).

Show this claim using a similar approach to our distinct elements analysis.

\end{document}
